\exposetitle{XIV}{Regular elements, continued. Application to algebraic groups}
\label{XIV}
\begin{center}by A. Grothendieck\\with an Appendix by J.-P. Serre\end{center}
{\renewcommand{\thefootnote}{0}\footnotetext{xy version of 1/12/08}}
\setcounter{footnote}{0}

% original p. 296
\sgasection{1}{Construction of Cartan subgroups and maximal tori for a smooth algebraic group}
\begin{theorem}{1.1}\label{XIV.1.1}
Let $G$ be a smooth algebraic group over a field $k$. Then $G$ admits a
maximal torus $T$, hence a Cartan subgroup $C=\Centr_G(T)$.
\end{theorem}
\begin{proof}
By XII 3.2, it amounts to the same thing to find a maximal torus $T$ of
$G$, or a Cartan subgroup $C$ of $G$. Moreover, since the maximal tori
of $G$ are those of $G^0$, one may suppose $G$ connected. We distinguish
two cases:

\noindent $1^\circ)$ The field $k$ is finite. Let $\mathcal T$ be the
scheme of maximal tori of $G$ (Exp XII 1.10.), which is a smooth scheme
over $k$. Note that $G$ acts on $\mathcal T$ by inner automorphisms,
and by the conjugacy theorem (XII 6.6 a)), two points of
$\mathcal T_{\bar k}$ rational over $\bar k$ are congruent under
$G_{\bar k}(\bar k)$. Taking into account that $\mathcal T$ is smooth
over $k$, hence $\mathcal T_{\bar k}$ smooth over $\bar k$, it follows
that $\mathcal T_{\bar k}$ is isomorphic to $G_{\bar k}/\bar N$, where
$\bar N$ is the stabilizer of an element of $\mathcal T_{\bar k}(\bar k)$,
i.e. the normalizer of a maximal torus $\bar T$ of $G_{\bar k}$.
Consequently, $\mathcal T$ is a ``homogeneous space'' under the action
of the group $G$. A well-known theorem of Lang (Amer. J. Math. 78,
1956, pp. 555--563) tells us that every homogeneous space under a
smooth connected algebraic group over a finite field $k$ admits a
rational point. In particular, $\mathcal T$ admits a rational point,
i.e. $G$ admits a maximal torus $T$. \hfill Q.E.D.

\noindent $2^\circ)$ The field $k$ is infinite. We shall use the
\begin{lemma}{1.2}\label{XIV.1.2}
Let $G$ be a smooth algebraic group over a field $k$. Then $G$ admits
a subgroup of type (C) (XIII 6.2).
\end{lemma}
By XIII 6.3 it amounts to the same thing to say that $\mathfrak g$
contains a Cartan subalgebra $\mathfrak d$. This is trivial if $k$ is
infinite, for then $\mathfrak g$ contains a regular element $a$, and
one will take $\mathfrak d=\operatorname{Nil}(a,\mathfrak g)$. The
case of finite $k$ is treated exactly as in proof $1^\circ)$ above,
but requires the preliminary construction of the scheme $\mathfrak d$
of Cartan subalgebras of $\mathfrak g$ and the fact that the latter
is smooth over $k$, which will be seen below (2.16). To establish
1.1 in case $2^\circ)$ in which we have placed ourselves, it will in
any event suffice to know 1.2 for infinite $k$. Let us also record:

\begin{lemma}{1.3}\label{XIV.1.3}
Let $G$ be a smooth connected affine algebraic group whose reductive
center (XII 4.1 and 4.4) is reduced to the unit group. Then $G$ is
nilpotent if (and only if) its Lie algebra $\mathfrak g$ is nilpotent.
\end{lemma}
This is contained in XII 4.9.

We can now give a method of constructing Cartan subgroups of $G$
(also valid when $k$ is finite, admitting 1.2 in this case). Suppose
first $G$ affine. We proceed by induction on $n=\dim G$, the assertion
being trivial if $n=0$. Thus suppose $n>0$ and the assertion proved
for dimensions $n'<n$. Let $Z$ be the reductive center of $G$, and let
\[
 u:G\longrightarrow G'=G/Z
\]
be the canonical homomorphism. By XII 4.7 c), one has a one-to-one
correspondence $C'\mto C=u^{-1}(C)$ between Cartan subgroups of $G'$
and Cartan subgroups of $G$. Thus, replacing $G$ by $G'$, one may
suppose that the reductive center $Z$ of $G$ is reduced to the identity
element (for this holds for that of $G'$ by XII 4.7 b)). By 1.3,
$G$ admits a subgroup $D$ of type (C). We know that over the algebraic
closure $\bar k$ of $k$, $D_{\bar k}$ contains a Cartan subgroup of
$G_{\bar k}$ (XIII 6.6 b)), so every Cartan subgroup of $D$ is a
Cartan subgroup of $G$ (XIII, 2.8 a)). Thus one is reduced to finding
a Cartan subgroup of $D$. If $\dim D=\dim G$, i.e. $D=G$, then the
Lie algebra of $G$ is a Cartan subalgebra of itself, hence is nilpotent,
so by 1.3 $G$ is nilpotent, hence is a Cartan subgroup of itself.
If $\dim D<\dim G$, then by the induction hypothesis there exists
a Cartan subgroup of $D$, which is therefore a Cartan subgroup of
$G$, which completes the proof in the case where $G$ is affine.
In the general case, let $Z$ be the center of $G$; then $G/Z=G'$
is affine (XII 6.1), and for every Cartan subgroup $C'$ of $G'$, its
inverse image $C$ in $G$ is a Cartan subgroup of $G$ (XII 6.6 e)).
One is reduced to finding a Cartan subgroup of the smooth affine
algebraic group $G'$, a case which has already been treated.
% typo?: kept as printed: u^{-1}(C), and reference 1.3 for existence of D.
\end{proof}

\begin{corollary}{1.4}\label{XIV.1.4}
Let $G$ be a smooth group scheme of finite type over an artinian scheme
$S$. Then $G$ admits a maximal torus $T$, hence a Cartan subgroup
$C=\Centr_G(T)$. Every torus $S$ in $G$ is contained in a maximal torus.
\end{corollary}
One may indeed suppose $S$ local with residue field $k$; then by 1.1,
$G_0=G\times_S\Spec(k)$ admits a maximal torus $T_0$; by IX 3.6 bis
and X 2.3, $T_0$ comes from a torus $T$ of $G$, which is clearly a
maximal torus. The last statement follows by applying the preceding
result to the centralizer of $S$, which is indeed smooth over $k$
(XI 2.4).
% typo?: kept as printed: S also names the torus, and smooth over k.

\begin{remarks}{1.5}\label{XIV.1.5}
\noindent a) We shall give below (3.20, 3.21 and XV) variants of 1.4
in the case where $S$ is not assumed artinian.

\noindent b) I do not know whether every algebraic group $G$
(not necessarily smooth) over a field $k$ admits a maximal torus.
The question only arises in characteristic $p>0$, and using 1.1 for
a smooth quotient group of the form $G'=G/I$, where $I$ is a suitable
radicial subgroup of $G$ (for example, the kernel of a sufficiently
high power of the Frobenius homomorphism), by taking the inverse image
in $G$ of a maximal torus of $G'$, one is reduced to the case where
$(G_{\bar k})_{\mathrm{red}}$ is a torus ($\bar k$ still denoting the
algebraic closure of $k$). It is easy to see that the answer is
affirmative when $G$ is commutative (or more generally nilpotent):
then $G$ admits a unique maximal torus, which one may construct,
for example, by descent from the maximal torus of $G_{\bar k}$%
{\renewcommand{\thefootnote}{*}\footnote{M. Raynaud has given a negative
answer to the question raised here; cf. XVII Example 5.9.c).}\addtocounter{footnote}{-1}}.

\noindent c) In the case where $G$ is affine, and $k$ perfect or $G$
solvable, 1.1 is known and due to Rosenlicht; his proof is very different
from that given here.

\noindent d) When $k$ is infinite, 1.1 is a consequence of the much
more precise result that the scheme $\mathcal T$ of maximal tori of
$G$ is a rational variety, proved below (6.1). The method is essentially
a combination of the proof of 1.1 and the explicit description of the
structure of the scheme $\mathfrak d$ of Cartan subalgebras of
$\mathfrak g$. To arrive at the desired result, we must first generalize
to the case of an arbitrary base prescheme certain results of XIII 4
to 6 (this is the purpose of the next two sections), and make the
preceding construction proving 1.1 more precise, using the fact that
every Cartan subgroup of $G$ is contained in one and only one subgroup
of type (C) of $G$ (Nos. 4, 5).
% typo: un en un seul -> un et un seul.
\end{remarks}

\sgasection{2}{Lie algebras over an arbitrary prescheme: regular sections and Cartan subalgebras}
\begin{lemma}{2.1}\label{XIV.2.1}
Let $A$ be a ring, $\mathfrak d$ a Lie algebra over $A$, and for every
$s\in\Spec(A)$, let $\mathfrak d(s)=\mathfrak d\otimes_A k(s)$ be the
corresponding Lie algebra over the residue field $k(s)$. One assumes
the $A$-module $\mathfrak d$ finitely presented. The following
conditions are equivalent:
\begin{enumeratei}
\item For every $s\in\Spec(A)$, $\mathfrak d(s)$ is nilpotent.
\item For every $x\in\mathfrak d$, $\ad(x)$ is nilpotent.
\item There exists an integer $N\geq0$ such that for every sequence
$x_1,\ldots,x_N$ of elements of $\mathfrak d$, one has
\[
 \ad(x_1)\ad(x_2)\ldots\ad(x_N)=0.
\]
\end{enumeratei}
\end{lemma}
When $A$ is a field, the equivalence of (i) and (iii) is the definition
of ``nilpotent'', and that of (ii) and (iii) is a well-known consequence
of Engel's theorem (Bourbaki, \emph{Groupes et alg\`ebres de Lie},
Chap. I, \S 4, \No 2). In the general case, one trivially has
(iii) $\Rightarrow$ (ii), and (ii) $\Rightarrow$ (i) by the preceding
result and the fact that (ii) is stable under passage to a quotient
and localization. It remains to prove (i) $\Rightarrow$ (iii).
When $A$ is artinian local with maximal ideal $\mathfrak m$, let
$n>0$ be an integer such that $\mathfrak m^n=0$, and $N$ an integer
such that condition (iii) holds for $\mathfrak d(s)=
\mathfrak d\otimes_A A/\mathfrak m$; take $N'=nN$; one sees at once
that this integer satisfies (iii). When $A$ is noetherian, there
then exist finitely many elements $s_i\in\Spec(A)$ and ideals of
definition $\mathfrak q_i\subset A_{s_i}$ such that the natural map
\[
 \mathfrak d\longrightarrow\prod_i\mathfrak d\otimes_A A_i,
 \qquad\text{where }A_i=A_{s_i}/\mathfrak q_i,
\]
is injective; then by what precedes there exists for every $i$ an
integer $N_i$ satisfying (iii) for the Lie algebra
$\mathfrak d\otimes_A A_i$, and taking for $N$ the largest of the
$N_i$, one satisfies (iii) for $\mathfrak d$. Finally, the general
case reduces to the noetherian case by the usual procedure
explained in EGA IV 8.
% typo: corps. l'equivalence -> corps, l'equivalence.

\begin{definition}{2.2}\label{XIV.2.2}
Let $S$ be a prescheme, $\mathfrak d$ a quasi-coherent Lie algebra
over $S$ which is a finitely presented module. One says that
$\mathfrak d$ is locally nilpotent if for every $s\in S$, the
Lie algebra $\mathfrak d(s)$ over $k(s)$ is nilpotent. One says
that $\mathfrak d$ is strictly locally nilpotent if it is locally
free, and if on every open subset $U$ of $S$ on which it has
constant rank $r$, its Killing polynomial reduces to
$P_{\mathfrak d}(t)=t^r$.
\end{definition}
Of course, if $\mathfrak d$ is a locally free Lie algebra (as a
module) over $S$, one defines its Killing polynomial as a polynomial
\[
 P_{\mathfrak d}\in A[t],\qquad
 A=\Gamma(\underline{\operatorname{Sym}}_{\OO_S}(\mathfrak d))
 \simeq\Gamma(\mathrm W(\mathfrak d))
\]
where $A$ is the ring of sections of the structure sheaf of the
vector bundle $\mathrm W(\mathfrak d)$ defined by $\mathfrak d$.
% typo?: kept as printed: Sym(d), without the dual.
It is clear that the two notions just introduced in 2.2 are stable
under base change and local for the fpqc topology. Note:

\begin{proposition}{2.3}\label{XIV.2.3}
If $\mathfrak d$ is strictly locally nilpotent, it is locally
nilpotent. The converse is true when $\mathfrak d$ is locally
free and $S$ reduced.
\end{proposition}
The proof is immediate.

\begin{definition}{2.4}\label{XIV.2.4}
Let $S$ be a prescheme, $\mathfrak g$ a Lie algebra over $S$
which is a locally free module of finite type. Let $\mathfrak d$
be a Lie subalgebra of $\mathfrak g$; one says that it is a
Cartan subalgebra if it satisfies the following conditions:
\begin{enumeratei}
\item $\mathfrak d$ is a submodule locally a direct summand
(hence also locally free of finite type).
\item For every $s\in S$, $\mathfrak d(s)$ is a Cartan subalgebra
of $\mathfrak g(s)$.
\end{enumeratei}
\end{definition}

\begin{definition}{2.5}\label{XIV.2.5}
Let $S,\mathfrak g$ be as in 2.4. A section $a$ of
$\mathfrak g$ is said to be quasi-regular if for every $s\in S$,
$a(s)\in\mathfrak g(s)$ is a regular element of the Lie algebra
$\mathfrak g(s)$ over $k(s)$. One says that $a$ is a regular
section if it is quasi-regular and contained in a Cartan
subalgebra of $\mathfrak g$.
\end{definition}
Notions 2.4 and 2.5 are again stable under base change and local
for the fpqc topology. Only the last assertion, and for the
case of the notion ``regular section'', requires a proof,
and follows from the fact that the Cartan subalgebra containing
a given regular section is uniquely determined. More precisely:

\begin{proposition}{2.6}\label{XIV.2.6}
Let $S,\mathfrak g$ be as in 2.4, and let $a$ be a quasi-regular
section of $\mathfrak g$. Then there exists at most one Cartan
subalgebra of $\mathfrak g$ containing $a$. For one to exist,
i.e. for $a$ to be a regular section, it is necessary and
sufficient that $a$ satisfy the following condition:
$\mathfrak d=\operatorname{Nil}(a,\mathfrak g)$ is a submodule
locally a direct summand of $\mathfrak g$, and $\ad(a)$ induces
an automorphism of $\mathfrak g/\mathfrak d$. In this case,
$\mathfrak d$ is the unique Cartan subalgebra of $\mathfrak g$
containing $a$.
\end{proposition}
Indeed, suppose that $a$ is contained in the Cartan subalgebra
$\mathfrak d$ of $\mathfrak g$. Then
$\ad(a)_{\mathfrak g/\mathfrak d}$ is bijective on every fiber,
hence (since $\mathfrak g/\mathfrak d$ is locally free of finite
type) it is an automorphism of $\mathfrak g/\mathfrak d$; on
the other hand, by 2.1, $\ad(a)_{\mathfrak d}$ is locally
nilpotent, hence $\mathfrak d=\operatorname{Nil}(a,\mathfrak g)$,
which proves the uniqueness of $\mathfrak d$ and the necessity
of the announced regularity criterion. For sufficiency, note
that the hypothesis made on $a$ implies that the formation
of $\operatorname{Nil}(a,\mathfrak g)$ commutes with extension
of the base and in particular with passage to the fibers,
which shows in particular that the fibers $\mathfrak d(s)$
of $\mathfrak d=\operatorname{Nil}(a,\mathfrak g)$ are Cartan
subalgebras of the $\mathfrak g(s)$; moreover $\mathfrak d$
is a Lie subalgebra of $\mathfrak g$ by XIII 4.1, hence it
is a Cartan subalgebra.

\begin{corollary}{2.7}\label{XIV.2.7}
Under the conditions of 2.4, let $\mathfrak d$ be a Cartan
subalgebra of $\mathfrak g$, $a$ a section of $\mathfrak d$
which is regular in $\mathfrak g$, and $u$ an automorphism
of $\mathfrak g$. For $u$ to leave $\mathfrak d$ invariant,
it is necessary and sufficient that $u(a)$ be a section
of $\mathfrak d$.
\end{corollary}
Indeed, then by transport of structure $u(a)$ is a regular
section of $\mathfrak g$, contained in both Cartan subalgebras
$\mathfrak d$ and $u(\mathfrak d)$, which are therefore identical.

\begin{corollary}{2.8}\label{XIV.2.8}
Under the conditions of 2.4, let $\mathfrak d$ be a Cartan
subalgebra of $\mathfrak g$. Then for every $s\in S$ such
that $k(s)$ is infinite, there exists an open neighborhood
$V$ of $s$ and a regular section $a$ of $\mathfrak g$ over
$V$ such that $\mathfrak d=\operatorname{Nil}(a,\mathfrak g|_V)$
(i.e. such that $a$ is a section of $\mathfrak d|_V$).
\end{corollary}
Indeed, the fact that $k(s)$ is infinite ensures the existence
of a regular element $a_0$ of $\mathfrak g(s)$ contained
in $\mathfrak d(s)$; one may extend it to a section $a$ of
$\mathfrak d$ over an open neighborhood $U$ of $s$, and
since $\ad(a)_{\mathfrak g/\mathfrak d}$ induces an
automorphism of $\mathfrak g(s)/\mathfrak d(s)$, it follows
that, provided one restricts $V$, $\ad(a)_{\mathfrak g/\mathfrak d}$
is an automorphism, which implies that $a$ is a quasi-regular
section of $\mathfrak g|_V$ satisfying the desired condition.

Let $\mathfrak g$ be as in 2.4; then examination of its
Killing polynomial implies at once that the function
\[
 s\mto\text{nilpotent rank of }\mathfrak g(s)
\]
on $S$ is upper semicontinuous. We are mainly interested
in the case where this function is in fact continuous,
i.e. locally constant. Here are some variants of this property:

\begin{proposition}{2.9}\label{XIV.2.9}
Let $S,\mathfrak g$ be as in 2.4. Consider the following conditions:
\begin{description}
\item[$(\mathrm C_0)$] The nilpotent rank of the $\mathfrak g(s)$
($s\in S$) is a locally constant function of $s$.
\item[$(\mathrm C_1)$] There exists a Cartan subalgebra of
$\mathfrak g$ locally for the fpqc topology.
\item[$(\mathrm C'_1)$] As $(\mathrm C_1)$, replacing ``fpqc
topology'' by ``\'etale topology''.
\item[$(\mathrm C_2)$] Condition $(\mathrm C_0)$ holds, and
for every $S'$ over $S$, every quasi-regular section of
$\mathfrak g_{S'}=\mathfrak g\otimes_S S'$ is regular.
\item[$(\mathrm C_3)$] Every $s\in S$ has an open neighborhood
$V$ on which the Killing polynomial of $\mathfrak g$ is of
the form
\[
 t^r(t^{n-r}+c_1t^{n-r-1}+\cdots+c_{n-r})
\]
where for every $s\in V$, $c_{n-r}(s)\in
\operatorname{Sym}(\mathfrak g(s)^\vee)$ is nonzero.
\end{description}
With this notation, one has the implications
\[
 (\mathrm C_3)\Rightarrow(\mathrm C_2)\Rightarrow(\mathrm C'_1)
 \Rightarrow(\mathrm C_1)\Rightarrow(\mathrm C_0),
\]
and when $S$ is reduced, the five conditions considered are equivalent.
\end{proposition}
\begin{proof}
Note also that the conditions considered are manifestly stable
under base change and local for the fpqc topology.

The implications $(\mathrm C'_1)\Rightarrow(\mathrm C_1)
\Rightarrow(\mathrm C_0)$ are trivial, and implication
$(\mathrm C_0)\Rightarrow(\mathrm C_3)$ when $S$ is reduced
is immediate. Note moreover:
\begin{corollary}{2.10}\label{XIV.2.10}
Suppose condition $(\mathrm C_0)$ holds. Let $U$ be the set
of elements of $\mathrm W(\mathfrak g)$ which are regular
in their fiber. Then $U$ is open; in particular, for every
section $a$ of $\mathfrak g$ over $S$, the set $V$ of
$s\in S$ such that $a(s)\in\mathfrak g(s)$ is regular is open.
\end{corollary}
Indeed, the first assertion follows from the second (applied
to $\mathfrak g_{S'}$ for every base change $S'$ over $S$).
For the latter, since one may here suppose $S$ reduced,
hence $(\mathrm C_3)$ verified, it suffices to examine the
Killing polynomial of $a$ in $\mathfrak g$.

Implication $(\mathrm C_2)\Rightarrow(\mathrm C'_1)$ now
follows easily from b) in the following more precise statement:
\begin{corollary}{2.11}\label{XIV.2.11}
Suppose condition $(\mathrm C_2)$ holds. Then:

\noindent a) For every $s\in S$ and every Cartan subalgebra
$\mathfrak d_0$ of $\mathfrak g(s)$ such that $\mathfrak d_0$
contains a regular element of $\mathfrak g(s)$ (a condition
automatically verified if $k(s)$ is infinite), there exists
an open neighborhood $V$ of $s$ and a subalgebra
$\mathfrak d$ of $\mathfrak g|_V$ whose fiber at $s$ is
$\mathfrak d_0$. If $S_1$ is a subprescheme of $S$ containing
$s$, and if one has already extended $\mathfrak d_0$ to
a Cartan subalgebra $\mathfrak d_1$ of
$\mathfrak g\otimes_S S_1$, then one may find an open
neighborhood $V$ of $s$ in $S$ and a Cartan subalgebra
$\mathfrak d$ of $\mathfrak g|_V$ such that
$\mathfrak d\otimes_V(S_1\cap V)$ is equal to
$\mathfrak d_1|_{(S_1\cap V)}$.
% typo: une sous-algebre de d -> une sous-algebre d.

\noindent b) For every $s\in S$ such that $\mathfrak g(s)$
contains a regular element (a condition automatically
satisfied if $k(s)$ is infinite), there exists an open
neighborhood $V$ of $s$ and a Cartan subalgebra
$\mathfrak d$ of $\mathfrak g|_V$.
\end{corollary}
Statement b) follows from a) by taking
$\mathfrak d_0=\operatorname{Nil}(a_0,\mathfrak g(s))$,
$a_0$ being a regular element of $\mathfrak g(s)$.
To prove a), say the second formulation, one considers
a regular element $a_0$ of $\mathfrak g(s)$ contained
in $\mathfrak d_0$, extends it in a neighborhood of
$s$ in $S_1$ to a section of $\mathfrak d_1$, and
extends the latter to a section of $\mathfrak g$ in
a neighborhood of $s$. By 2.10, this section is
quasi-regular in an open neighborhood $V$ of $s$,
hence regular by $(\mathrm C_2)$, so by 2.6,
$\mathfrak d=\operatorname{Nil}(a,\mathfrak g|_V)$
satisfies the desired conditions.
% typo: section et quasi-reguliere -> section est quasi-reguliere.
% Proof of Proposition 2.9 continues in en-02.tex.
